\section{Related Work}
\subsection{From Latent Reasoning to Concept-Level Language Modeling}

Recent research has explored performing reasoning at higher levels of abstraction than individual tokens, offering both computational efficiency and new modeling capabilities. \textbf{Latent reasoning} frameworks perform reasoning entirely within continuous hidden state spaces rather than through explicit token generation~\cite{hao2025traininglargelanguagemodels}. In the COCONUT framework, the model's hidden state from one reasoning step feeds directly into subsequent steps without generating intermediate tokens~\cite{hao2025traininglargelanguagemodels}. This approach offers significant computational advantages over methods like Chain-of-Thought prompting, which require generating hundreds of intermediate tokens. More importantly, continuous representations can encode multiple potential reasoning paths in superposition, enabling parallel exploration of the solution space~\cite{hao2025traininglargelanguagemodels}. However, this comes at the cost of interpretability and can struggle with tasks requiring precise symbolic manipulation.

Building upon latent reasoning principles, the \textbf{Large Concept Model (LCM)} framework operates at an intermediate level: reasoning on sentence-level "concepts" rather than tokens~\cite{lcm2024large}. LCMs use a three-stage pipeline: (1) a frozen encoder maps sentences to fixed-size embeddings in a semantic space (e.g., SONAR, supporting 200 languages), (2) a transformer performs autoregressive prediction in this concept space using diffusion or quantization techniques adapted from computer vision, and (3) a frozen decoder reconstructs token-level output~\cite{lcm2024large}. This approach combines key advantages: like latent reasoning, it operates in continuous space with substantial efficiency gains (10× sequence length reduction), but unlike pure latent reasoning, each concept remains interpretable as a decodable sentence~\cite{lcm2024large}. Remarkably, LCMs demonstrate zero-shot multilingual transfer—models trained only on English can generate in 200+ languages by leveraging the language-agnostic semantic space~\cite{lcm2024large}. However, the LCM framework faces significant limitations. First, it requires pre-training separate encoder and decoder models on massive multilingual data before the LCM itself can be trained, creating a scalability bottleneck~\cite{lcm2024large}. Second, and more fundamentally, the sentence-level granularity is a fixed human prior—the model must accept predetermined sentence boundaries rather than learning task-optimal segmentation. This rigidity prevents the model from adapting its conceptual granularity to different domains or tasks. Our DCLM architecture addresses both limitations through end-to-end training with dynamic, learnable boundary detection that discovers optimal chunking strategies directly from the data.

\subsection{Dynamic Compute Allocation in Language Models}

Standard LLMs allocate uniform computation to every token, ignoring that some tokens (e.g., predictable function words) require minimal processing while others (e.g., concept boundaries) demand more effort~\cite{kaplan2020scaling, hoffmann2022training}. Recent work explores adaptive allocation mechanisms. The Universal Transformer~\cite{dehghani2018universal} introduced recurrence in depth, applying the same transformation block repeatedly with a learned halting mechanism that determines when each position has been sufficiently refined. Mixture of Experts (MoE) models~\cite{shazeer2017outrageously, jiang2024mixtral} achieve conditional computation by routing each token to a subset of expert sub-networks. However, these approaches focus on parameter efficiency and scaling rather than fundamentally addressing the information density problem.

\textbf{H-NET}~\cite{hwang2025dynamicchunkingendtoendhierarchical} directly addresses adaptive allocation through learned boundary detection. The model predicts semantic boundaries by analyzing local patterns (e.g., similarity between consecutive hidden states), segments the sequence into variable-length chunks, and processes the compressed chunk representations hierarchically. Critically, boundary detection is differentiable and trained end-to-end, allowing task-appropriate chunking strategies to emerge~\cite{hwang2025dynamicchunkingendtoendhierarchical}. This yields substantial gains: learned boundaries align with linguistic structures even without supervision, and compression (4-8× reduction) translates to quadratic attention savings~\cite{hwang2025dynamicchunkingendtoendhierarchical}. The approach implicitly allocates more computation to high-information boundaries where new concepts begin—precisely where prediction is most difficult. However, H-NET's primary focus is on efficient representation through hierarchical bit-level modeling rather than token-level generation in state-of-the-art autoregressive LLMs~\cite{hwang2025dynamicchunkingendtoendhierarchical}. This leaves unaddressed the critical problem of computational waste in modern decoder-only language models, where every token—regardless of its predictability or information content—receives identical processing through the full model depth. Our DCLM architecture bridges this gap by adapting H-NET's dynamic boundary detection principles to the token-level generation paradigm of current LLMs~\cite{brown2020language, touvron2023llama}. By segmenting sequences into concept chunks and performing compressed reasoning before token-level decoding, DCLM enables adaptive computation allocation in the exact architectural context where it matters most: next-token prediction in large-scale autoregressive models.